\section{The Beautiful Losers of History}

\begin{quotex}
Then know that Allah has described Himself as the Outwardly Manifest and the Inwardly Hidden; He brought the universe into existence as a Visible world and an Unseen world so that we might know the Hidden by the Unseen and the Manifest by the Visible. \flright{\textsc{Ibn Arabi}, \emph{The Wisdom of the Prophets}}

\end{quotex}
\paragraph{Schopenhauer}
Much is missed if you don't read the deepest thinkers directly. You miss the motivation, the problems that need to be solved, and so on. One image from Schopenhauer's masterwork remains with me. He relates the story of a man in Morocco whose lower jaw was pulled from his face. I don't know how long news like that took to reach Germany, but how, just how, does that event fit into a philosophical system? That is what he tried to answer.

\subparagraph{The World as Representation}
If a tree falls in the forest, and no one is nearby to hear it, does it make a sound? You may have already answered that, although it produces sound waves, there is no sound.

If a tree falls in the forest, and no one is nearby to see it, does it make a sight? The answer ``no" will elicit doubt from most people.

Yet if the former is true, then so must be the latter. Following Spinoza, if Extension has no qualities, then qualities must be a product of Thought. This does not imply that the world is imaginary or arbitrary. Analogous to Jung's Collective Unconscious, there is a Collective Conscious that experiences the world.

This is the outwardly manifest world.

Kant called that the phenomenal world. He speculated that there is an unseen world — the noumenal world — that corresponds one-to-one with the phenomenal world. That makes no sense. If objects are creations of the Mind — and Extension has no qualities — postulating invisible objects corresponding to manifest objects make no sense. The world represented in the Mind simply is that world.

This, Schopenhauer asserts, must be grasped intuitively as a Principle. He writes:

\begin{quotex}
``The world is my representation" is, like Euclid's axioms, a proposition which everyone must recognize as true as soon as he understands it, although it is not a proposition that everyone understands as soon as he hears it.

\end{quotex}
\subparagraph{Space, Time, Causality, Will}
Space and time are, from the metaphysical perspective, preconditions for the human state.

Space, time, and causality are likewise mental phenomena. That is, there is no space independent of mind that gets filled up with stuff, nor time without any objects. Things relate to each other in terms of space, time, or cause. It is unnecessary to resort to ``matter" as the causal element for the objects of our experience.

For Descartes, the primary object of our experience is ``I am". Schopenhauer takes it once step further, identifying our Will as our primal experience. This is consistent with Spinoza's notion of ``conatus", i.e., the striving of things to preserve themselves in being. This is experienced in human beings as the conscious experience of desire, human freedom, and good and evil. Thus the human being strives for an increase in power, which is the same as virtue. To the extent that he recognizes the source of his motives, he gains power over himself.

Thought occurs successively in time, so it is how the subject comes to know itself. The will is identical in time and manifests itself over the course of life. He asserts:

\begin{quotex}
The course of life retains the same fundamental tone; in fact, its manifold events and scenes are at bottom like variations on one and the same theme.

\end{quotex}
Thus life is like a symphony or an opera, in which themes reappear in slightly different forms. For self-understanding, it is necessary to pay attention to those themes. And like a good opera, the main theme ultimately is resolved. You can consciously act out your role or you can just let it happen to you.

Our will and our acts are not two different things, as some assume that the will is the ``cause" of the movements of the body. The Will is to the Act, as the inside is to the outside. The gnomic will, the feeling of deliberation of alternatives, often accompanied by anxiety or inertia, is not at all our True Will. Our true will is the spontaneous expression of who we are. When the proper circumstances appear that are in conformance with one's destiny, then the True Will seizes the opportunity. The gnomic will delays and the opportunity is lost.

\subparagraph{Schopenhauer's God}
So if things, space, time, cause, and even one's own self are aspects of the Manifest World, Schopenhauer turns to what is Hidden, or the Noumenal, as he calls it. Matter is part of the phenomenal world. But matter is in motion everywhere we look, from the stars and planets down to animals, plants, and so on; hence, it can't be ultimate. Rather it is energy that is ultimate in the phenomenal world. Matter, therefore, is just condensed energy, that is, energy confined to a particular space.

Now the only thing that we are aware of, from the inside, is our own will. It is clear that we don't know what it is like to be a ``bat", for example, or a stone, etc. Hence, the conclusion is that ultimate reality must be related to the Will. So energy as phenomena is condensed Will, as noumenon; this Will is not the same as the human will. Rather, it is a blind urge, purposeless, a striving devoid of knowledge.

As it is Hidden and Noumenal, we can have no direct experience of the Will. Rather, our experience of it consists of the random motions and events of the phenomenal world. It is easy to see our world as random, cruel, and without any particular goal.

Just like Spinoza's God, Schopenhauer's Will is a challenge to Traditional teaching since they mimic it in many respects. The difference is that the philosophers believe that they can comprehend the essence of the Absolute, or the Unseen world, through reason alone.

\subparagraph{World Denial as Way of Life}
In the face of the hopelessness of the blind will, Schopenhauer retreats into a profound pessimism. His only moral principle is to cause no harm, to bring no more suffering into the world. The obviously corollary to that is his antinatalism. Life brings suffering. Although he admired Buddhism for its pessimistic outlook on life, he failed to accept its teaching on the end of suffering. Besides Buddhism, he also found solace in the mysticism of Brahmanism and also of neoplatonic currents in Christianity.

His escape valve was his love of the arts and High Culture, particularly music. That is a countercurrent, a swimming against the tide, of the force of the Will.

\subparagraph{Power as Virtue}
In our time, particularly in the Western countries, a philosophy akin to Schopenhauer's, is dominant, usually unconsciously. Thus birth rates have plummeted. However, there is no longer High Culture as solace, since the blind will is even destroying that.

There are many naïve souls who believe that logic, evidence, facts, persuasion, objectivity, and the like, will convince the Will to behave. Only power can oppose power. Until the revolters against the modern world recognize it, they will be counted amongst the beautiful losers of history.



\flrightit{Posted on 2021-04-18 by Cologero }

\begin{center}* * *\end{center}

\begin{footnotesize}\begin{sffamily}



\texttt{Greg on 2021-04-21 at 12:11 said: }

Only power can oppose power. A sobering truth. It appears that all power resides with the elites who have shaped the modern world and will continue dragging it down. So we are destined to be beautiful lovers? What of it? What good is it too be on the right side of history if that places you on the wrongside of eternity?Apocalyptic literature suggests that history ends with the triumph of darkness and evil. It may be that to choose God, one must surrender the world. Or perhaps you mean supernatural power, sacred magic as Tomberg calls it. Perhaps a battle between sacred power and the powers of the world?


\hfill

\texttt{Cologero on 2021-04-29 at 07:44 said: }

Who knows, Greg. Morality seems to be a luxury, suitable only for life in a civilization. Otherwise, raw power seems to hold sway.


\hfill

\texttt{Angolmois on 2021-05-02 at 14:10 said: }

Greg, there was a line in Robert Bolton's The Order of the Ages that caught my eye and which went something like this: ``it is no great feat to be spiritual in a spiritual civilization where everything anti-spiritual is suppressed; the true way is found only when almost everything is opposed to it." Spiritual people must be living the best times since at least one's convictions are being challenged from every possible direction.


\hfill

\texttt{Jorgen B on 2021-05-23 at 23:31 said: }

Schopenhauer would have been a monk in the middle ages.


\hfill

\texttt{Reza on 2021-06-21 at 22:51 said: }

@Jorgen

Not so my friend! Don't confuse his philosophy with his personality and inclinations . From his biographies one comes to know the man to be a bon vivant.


\end{sffamily}\end{footnotesize}
